% Гармонический осциллятор: слева -- траектории в фазовой плоскости (эллипсы),
% справа -- интегральная кривая в расширенном фазовом пространстве (винтовая
% линия над эллипсом). Ответ на вопрос в конце лекции.
%
% Правая панель нарисована чистым TikZ с заданными базисными векторами x/y/z:
% addplot3 из pgfplots на такой кривой упирается в main memory size.
\begin{figure}[htbp]
\centering
\begin{tikzpicture}[>={Latex[length=4pt]}]
\begin{axis}[
  width=6.2cm, height=5.4cm,
  axis lines=middle,
  xmin=-2.3, xmax=2.3, ymin=-2.5, ymax=2.5,
  xlabel={$x$}, ylabel={$v$},
  xlabel style={anchor=north west}, ylabel style={anchor=south east},
  xtick=\empty, ytick=\empty, clip=false,
  title style={font=\footnotesize, yshift=-4pt},
  title={(а) фазовая плоскость $\R^2$},
]
  \foreach \r in {0.6, 1.1, 1.6} {
    \addplot[thick, blue!60!black, domain=0:360, samples=90]
      ({\r*cos(x)}, {1.3*\r*sin(x)});
  }
  \addplot[only marks, mark=*, mark size=1.2pt, black] coordinates {(0,0)};
  \draw[->, red!65!black, thick] (axis cs:1.1,0.30) -- (axis cs:1.04,0.75);
\end{axis}
\end{tikzpicture}%
\hfill
\begin{tikzpicture}[
  >={Latex[length=4pt]},
  x={(0.80cm,-0.30cm)}, y={(0.62cm,0.34cm)}, z={(0cm,0.62cm)},
]
  % оси
  \draw[->] (0,0,0) -- (2.6,0,0) node[anchor=north, font=\footnotesize] {$x$};
  \draw[->] (0,0,0) -- (0,2.6,0) node[anchor=west, font=\footnotesize] {$v$};
  \draw[->] (0,0,0) -- (0,0,6.6) node[anchor=south, font=\footnotesize] {$t$};

  % траектория: проекция на фазовую плоскость
  \draw[gray, dashed, thick]
    plot[variable=\s, domain=0:360, samples=72, smooth cycle]
    ({1.5*cos(\s)}, {1.5*sin(\s)}, 0);

  % интегральная кривая: винтовая линия над эллипсом
  \draw[very thick, blue!60!black]
    plot[variable=\s, domain=0:720, samples=145, smooth]
    ({1.5*cos(\s)}, {1.5*sin(\s)}, {\s/720*5.8});

  \node[gray, anchor=north west, font=\footnotesize] at (1.7,-1.5,0)
    {траектория};
  \node[blue!60!black, anchor=west, font=\footnotesize] at (1.6,0.6,5.6)
    {интегральная кривая};
  \node[anchor=south, font=\footnotesize] at (0.6,0,7.4)
    {(б) расширенное фазовое пространство $\R^3$};
\end{tikzpicture}
\caption{Гармонический осциллятор. (а) Траектории в фазовой плоскости —
эллипсы $\omega^2 x^2 + v^2 = \alpha^2\omega^2$; система автономна, поэтому
траектория не зависит от начального момента $t_0$. (б) Интегральная кривая в
расширенном фазовом пространстве — винтовая линия над этим эллипсом; сдвиг
начального момента сдвигает её вдоль оси $t$, не меняя проекции.}
\label{fig:oscillator-phase}
\end{figure}
